\section{RQ3: Does the anchored model survive the same tests?}
\label{sec:stress}

If the diagnosis is right, anchoring to something the transform cannot change should move both
tests specifically where the input has left the training distribution.

% >>> GENERATED TABLE rq4_ladder -- from tables/rq4_ladder.tex; regenerate, then re-run inline_tables.py
% GENERATED 2026-09-14 02:18 UTC by scripts/paper/gen_fse_tables.py
% SOURCE: results/analysis/pipeline/{f2_divergence_codellama7b,composite_depth_codellama7b}.json
% Do not edit by hand -- regenerate. New: the stimulus this table describes did not exist when the draft was written.
\begin{table*}[ht]
\centering\footnotesize\setlength{\tabcolsep}{3pt}
\caption{\textbf{RQ4: the divergence ladder.} Composites that contain a transform family the model
has never seen. Breadth's advantage flips sign the moment such a component enters, while anchoring
rises \emph{above} the clean-code control. The pre-registered rule required both that anchoring hold
and that breadth fall below the control; only the first is established, since the pooled interval at
d2 straddles zero. CodeLlama-7B, 287 programs common to all six composites, fixed before any system
was evaluated.}
\label{tab:rq4_ladder}
\begin{tabular}{@{}llccc@{}}
\toprule
\textbf{Level} & \textbf{Stimulus} & \texttt{mono}$-$\texttt{clean} & \texttt{cons}$-$\texttt{clean} & \texttt{cons}$-$\texttt{mono} \\
\midrule
d0 & depth 2, all seen & \textbf{$+3.49$ [$+2.04$, $+5.01$]} & -- & $+0.87$ [$-0.59$, $+2.33$] \\
d1 & depth 3--4, all seen & \textbf{$+4.15$ [$+2.37$, $+5.91$]} & \textbf{$+5.02$ [$+3.41$, $+6.69$]} & $+0.87$ [$-0.59$, $+2.33$] \\
\midrule
d2 & depth 2, \textbf{unseen inside} & $-1.71$ [$-3.60$, $+0.31$] & \textbf{$+2.13$ [$+0.66$, $+3.60$]} & \textbf{$+3.84$ [$+1.90$, $+5.70$]} \\
d3 & depth 3, \textbf{unseen inside} & $-0.35$ [$-2.67$, $+1.86$] & \textbf{$+3.26$ [$+0.81$, $+5.92$]} & \textbf{$+3.60$ [$+1.40$, $+5.81$]} \\
\bottomrule
\end{tabular}
\end{table*}
% <<< GENERATED TABLE rq4_ladder

\textbf{Composition.} Table~\ref{tab:rq4_ladder} is the divergence ladder, from all-seen depth-2
stacks (d0) to stacks containing a family the model has never seen (d2, d3); the d2 and d3 stimulus
did not exist before this work. Once an unseen component enters, breadth's advantage over the
clean-code control \textbf{flips sign}, from $+3.49$ and $+4.15$ to $-1.71$ and $-0.35$, and
anchoring, which merely matches breadth on seen stacks ($+0.87$ [$-0.59$, $+2.33$]), rises
\textbf{above the control} ($+2.13$, $+3.26$) and \textbf{above breadth} ($+3.84$, $+3.60$). The
pre-registered rule required both that anchoring hold and that breadth fall below the control; only
the first is established, because the pooled d2 interval spans zero, so by that rule the outcome is
\textsc{undecided}. Split by RQ1's identifier rule, committed before the first d2 cell was written,
identifier-containing stacks sit at $+1.98$ [$-0.12$, $+4.18$] and structural-containing stacks at
$-3.88$ [$-6.08$, $-1.71$]: breadth is broken where the stack leaves it no identifier surface to key
on. Across the panel (Table~\ref{tab:rq4_by_model}, Appendix~\ref{app:supp}) the anchoring advantage on unseen-containing
stacks replicates on five of eight models, is inconclusive on three and refuted on none; family
exposure is significant on all eight and is the largest effect in the table.


\textbf{Reversal.} In Table~\ref{tab:direction_ratio}, anchoring's direction ratio is $+0.03$: the
only arm in the supervised-fine-tuning family whose forward gain cost nothing backwards. It does
\textbf{not} acquire backward competence ($+0.62$ [$-1.56$, $+2.81$] against the untuned model);
the claim is that every existing method pays backward competence for its forward gain and anchoring
is the first that does not. The one arm that gains in both directions is trained on the unseen
family's sibling ($+1.63$ [$+0.11$, $+3.20$] backwards): exposure to a family transfers in a way
exposure to transforms does not.

\textbf{The bound.} On the unseen family alone the anchored model is as robust as clean-code tuning
and no more ($-0.30$, $-0.33$, $+0.74$ against the control at 7B, 13B, 34B; seventh of thirty-three
arms), and every arm that beats it there was trained on the family's sibling. Anchoring buys
breadth's recombination benefit at little cost, extends it to composed inputs the model has never
seen, and stops paying the backward price --- and it does not add the capability breadth lacks. It
is a repair of what adaptation breaks, not a route to invariance.

\begin{takeaway}{Takeaway: RQ3}
The anchored model's advantage appears only where the input diverges from training: level with breadth
on seen stacks, above the clean-code control and above breadth once an unseen family enters (five of
eight models replicate, none refute; the strict pre-registered rule is undecided), and a direction
ratio of $+0.03$ where every other supervised arm is negative. On the unseen family alone it is as
robust as clean-code tuning and no more: a repair of what adaptation breaks, not a route to invariance.
\end{takeaway}
